\section{Ergebnisse \& Evaluation}\label{sec:ergebnisse}

\subsection{Was funktioniert, was nicht?}

\subsubsection{Was funktioniert}

Die Implementierung zeigt mehrere erfolgreiche Aspekte.
Der \textbf{Agent vs Human Modus} ermöglicht eine erfolgreiche Integration von trainierten Strategien aus der Bachelorarbeit.
Strategien werden korrekt geladen und der Agent spielt automatisch basierend auf der trainierten Strategie.
Die zufällige Zuweisung von Player-IDs sorgt dafür, dass der menschliche Spieler nicht immer derselbe Spieler ist, was die Spielerfahrung verbessert.
Die Integration funktioniert nahtlos, da die Strategien direkt aus der Bachelorarbeit wiederverwendet werden können, ohne dass eine Konvertierung erforderlich ist.

Der \textbf{Human vs Human Modus} funktioniert zuverlässig, wobei der WebSocket-Server Echtzeit-Multiplayer-Spiele ermöglicht.
Zwei Clients können sich verbinden und ohne merkbare Latenz gegeneinander spielen.
Die Push-basierte Kommunikation funktioniert wie erwartet, und State-Updates werden sofort an beide Clients gesendet.
Die automatische Wiederverbindung bei Verbindungsfehlern sorgt für eine robuste Verbindung.

Die \textbf{Wiederverwendung} der Game-Logik aus der Bachelorarbeit funktioniert ohne Änderungen, was die modulare Architektur bestätigt.
Die Game-Logik wurde erfolgreich in eine separate Klasse (\texttt{PokerGameLogic}) gekapselt, die sowohl von HTTP- als auch von WebSocket-Servern verwendet werden kann.
Dies ermöglicht Code-Wiederverwendung und vereinfacht die Wartung erheblich.

Die \textbf{Tests} bieten umfassende Test-Abdeckung für Server und Game-Logik, wobei alle automatisierten Tests bestehen.
Die szenarien-basierten Tests sind verständlich und geben guten Aufschluss über die Funktionalität des Systems.
Die Tests können ohne GUI-Fenster ausgeführt werden, was die Integration in CI/CD-Pipelines ermöglicht.

Die \textbf{Thread-Safety} wurde erfolgreich implementiert, wobei keine Race Conditions bei Multiplayer-Server beobachtet wurden.
\texttt{threading.Lock} verhindert erfolgreich gleichzeitige Zugriffe auf kritische Datenstrukturen.
Die Lock-Hold-Zeit ist kurz, da kritische Abschnitte auf das Minimum beschränkt wurden, was den Performance-Overhead minimiert.

\textbf{Sicherheitsfeatures} wie Input-Validierung und Rate Limiting funktionieren wie erwartet.
Input-Validierung verhindert erfolgreich XSS-Angriffe durch HTML-Escaping, und Rate Limiting schützt vor Missbrauch durch zu viele Requests.

\subsubsection{Was nicht funktioniert / Einschränkungen}

Es gibt jedoch auch Einschränkungen.
Der Multiplayer-Modus unterstützt aktuell nur Heads-Up (2 Spieler).
Eine Erweiterung auf mehr Spieler würde Änderungen an der Game-Logik erfordern, da die aktuelle Implementierung für zwei Spieler optimiert ist.
Die Game-Logik aus der Bachelorarbeit unterstützt ebenfalls nur zwei Spieler, sodass eine Erweiterung auf mehr Spieler auch Änderungen an der Game-Logik erfordern würde.

Spielstände werden nicht gespeichert, sodass nach einem Neustart des Servers alle Spielzustände verloren gehen.
Dies ist eine bewusste Design-Entscheidung für diese Version, da Persistenz zusätzliche Komplexität erfordern würde, wie Datenbank-Integration oder Datei-System-Persistenz.

Es gibt keine Replay-Funktionalität, sodass Spiele nicht nachträglich angeschaut werden können.
Dies wäre eine nützliche Funktion für Analyse und Lernen, würde jedoch zusätzliche Implementierung erfordern, wie die Aufzeichnung von Spielzuständen und eine Wiedergabe-Funktionalität.

Große Strategie-Dateien benötigen Zeit zum Laden, was besonders bei Strategien mit vielen Iterationen spürbar ist.
Die Kompression hilft zwar, die Dateigröße zu reduzieren, aber das Laden und Dekomprimieren großer Dateien kann einige Sekunden dauern.
Dies ist jedoch akzeptabel, da Strategien nur einmal beim Start geladen werden.

Der Server hat keine Authentifizierung implementiert, sodass jeder sich verbinden kann, solange der Server nicht voll ist.
Dies ist für lokale Tests ausreichend, würde jedoch für Produktions-Einsatz zusätzliche Sicherheitsfeatures erfordern.

\subsection{Messungen, Tests und Beobachtungen}

\subsubsection{Performance-Beobachtungen}

Komprimierte \texttt{.pkl.gz} Dateien reduzieren den Speicherbedarf erheblich.
Eine unkomprimierte Strategie-Datei kann mehrere hundert MB groß sein, während die komprimierte Version oft unter 50 MB liegt.
Die Ladezeit bleibt akzeptabel (wenige Sekunden), da Strategien nur einmal beim Start geladen werden.
Die Kombination aus Pickle und gzip ermöglicht effiziente Speicherung und schnelles Laden großer Strategie-Dateien.

WebSocket zeigt deutlich niedrigere Latenz als HTTP-Polling.
Bei HTTP-Polling mit 100ms Intervall ist die Latenz mindestens 100ms, da Updates nur beim nächsten Polling-Intervall empfangen werden.
WebSocket hingegen ermöglicht sofortige Übertragung von Updates, was die Latenz praktisch eliminiert.
Dies verbessert die Spielerfahrung erheblich, da State-Updates sofort übertragen werden.

\texttt{threading.Lock} verhindert Race Conditions effektiv und hat minimalen Performance-Overhead.
Die Lock-Hold-Zeit ist kurz, da kritische Abschnitte auf das Minimum beschränkt wurden.
Selbst bei vielen gleichzeitigen Requests wurde keine merkbare Verlangsamung beobachtet.
Dies bestätigt, dass \texttt{threading.Lock} eine effektive Lösung für Thread-Safety in diesem Anwendungsfall ist.

Die GUI bleibt auch während Netzwerk-Operationen responsiv, da \texttt{QThread} für WebSocket-Kommunikation verwendet wird und die GUI-Thread nicht blockiert wird.
Dies ermöglicht eine flüssige Benutzererfahrung, auch wenn Netzwerk-Operationen im Hintergrund stattfinden.

\subsubsection{Test-Ergebnisse}

Alle automatisierten Tests bestehen zuverlässig.
Während der Testphase wurden keine Bugs in der Server-Logik gefunden.
Die Input-Validierung und Rate Limiting funktionieren wie erwartet.
Die HTML-Escaping-Funktionalität verhindert erfolgreich XSS-Angriffe.
Tests mit bösartigen Eingaben (z.B. \texttt{<script>alert('XSS')</script>}) werden korrekt escaped und stellen keine Sicherheitsrisiko dar.

Rate Limiting funktioniert korrekt.
Clients, die zu viele Requests senden, werden abgelehnt, während normale Nutzung nicht beeinträchtigt wird.
Dies schützt den Server vor Missbrauch und stellt sicher, dass die Ressourcen fair verteilt werden.

Die Tests decken die wichtigsten Szenarien ab: Client-Registrierung, State-Updates, Action-Verarbeitung, Reset-Funktionalität, Fehlerbehandlung und Sicherheitsfeatures.
Die szenarien-basierten Tests sind verständlich und geben guten Aufschluss über die Funktionalität des Systems.
